% ---------------------------------------------------------------------
%  03-hpb-in-one-page.tex -- the whole method on one page
% ---------------------------------------------------------------------
\section{HPB in One Page}
\label{sec:onepage}

\paragraph{The idea.}
Draw the product, not the work. Each product reference is a thick
horizontal line, the \emph{main line}, running left to right in time. Its
versions sit on it as \emph{baselines}: released ones on the left of the
NOW marker, planned ones on the right. Everything you do to the product is
a \emph{branch}: it leaves the main line at the released baseline it
started from and comes back to the version that closes it, or ends in a
cross with a written reason. Improvements and lab work run above the line,
defects and tasks below. When a hardware change creates a new reference,
the line itself forks and the new reference starts its own main line,
linked to its parent. A planned version is a \emph{gate}: it is released
only when the work committed to it has passed verification.

That is the whole method. \Cref{fig:anatomy} shows every element on a
single product line, ``Product A''. Read it from left to right: v1.0 is
released; defect X-31 forks from it and is fixed in v1.0.1; improvement
I-12 forks from v1.0.1 and is merged into v1.1; task T-7 is discarded
because it was superseded; lab evaluation L-3 is still open, with one test
passed and one failed; and defect X-40 is committed to the planned v1.2,
whose gate is still closed.

% F1 -- Anatomy of an HPB diagram: the glyphs of a single-line diagram.
% Illustrative data. Requires hpb.sty v1.1. Canvas: document defaults.
% Callout leaders never cross a label or a branch other than their target.
\begin{figure}[t]
\centering
\hpbfit{%
\begin{hpbgraph}
  % ---- lanes: proactive work above, reactive work below the main line
  \hpblane{2}{lab}{Lab (L)}{0}{19.6}
  \hpblane{1}{imp}{Improvement (I)}{0}{19.6}
  \hpblane{-1}{iss}{Issue (X)}{0}{19.6}
  \hpblane{-2}{task}{Task (T)}{0}{19.6}
  % ---- main line: history (solid) up to NOW, plan (dashed) afterwards
  \hpbmainline{0}{16}
  \hpbmainline*{16}{19.6}
  \hpborigin{0}{Product A}
  \hpbrelease{1}{v1.0}   \hpbcollapsed{1}{12}
  \hpbrelease{4}{v1.0.1}
  \hpbrelease{11}{v1.1}
  \hpbplanned{18}{v1.2}{Mar 2027}
  % ---- branches (upward forks from v1.0.1 clear its wide tag)
  \hpbbranch[label={X-31 Crash}]{iss}{1}{-1}{2}{3}{merge}{4}
  \hpbbranch[forkoff=0.5, label={T-7 Migrate logs}]{task}{4}{-2}{5}{14}{discard}{}
  \hpbbranch[forkoff=0.95, mergeoff=0.6, label={I-12 Faster pairing}]{imp}{4}{1}{6}{10}{merge}{11}
  \hpbbranch[forkoff=0.7, label={L-3 New sensor}]{lab}{4}{2}{7}{16}{open}{}
  \hpbbranch[label={X-40 Battery drain}]{iss}{11}{-1}{12}{16}{target}{18}
  % ---- steps on branches
  \hpbstep{lab}{8}{2}{pass}
  \hpbstep{task}{9}{-2}{done}
  \hpbstep{lab}{13}{2}{fail}
  \hpbstep{iss}{15}{-1}{todo}
  % ---- rationale of the discard, gate of the planned baseline, NOW
  \hpblabel[below center]{14}{-2}{\itshape rationale: superseded}
  \hpbgate[1.6]{18}{1}{-1}{gate v1.2}
  \hpbnow{16}{2}{-2}
  % ---- callouts, upper row (targets reachable without crossing a path)
  \hpbcallout{-1.15,3.35}{lane legend}{-1.15,2.3}
  \hpbcallout{1.25,3.35}{main line\\(the product)}{0.45,0.05}
  \hpbcallout{3.6,3.35}{version tag\\(released)}{1.15,0.95}
  \hpbcallout{6.2,3.35}{fork from\\baseline}{4.7,1.3}
  \hpbcallout{11.0,3.35}{lab test steps\\pass / fail}{13,2.15}
  \hpbcallout{14.3,3.35}{open branch\\(head at NOW)}{16,2.2}
  \hpbcallout{18.0,3.35}{planned baseline\\and target date}{18,1.0}
  % ---- callouts, lower rows
  \hpbcallout{-1.15,-3.25}{lane guide}{-0.05,-2}
  \hpbcallout{2.6,-3.25}{start node\\with type letter}{2,-1.2}
  \hpbcallout{7.6,-3.25}{sub-task step\\(done)}{9,-2.15}
  \hpbcallout{11.4,-3.25}{discard with\\rationale}{13.1,-2.45}
  \hpbcallout{16.3,-3.25}{pending step\\(todo)}{15,-1.15}
  \hpbcallout{0.4,-4.15}{collapsed history:\\12 merged branches}{1,-0.6}
  \hpbcallout{4.6,-4.15}{merge into\\baseline}{3.5,-0.6}
  \hpbcallout{18.45,-4.15}{targeted branch\\(dashed to gate)}{17.5,-0.6}
\end{hpbgraph}}
\caption{Anatomy of an HPB diagram, on one illustrative product line. The
thick line is the product and its versions; improvements and lab work run
above it, issues and tasks below. Every branch starts at the version it
forked from and ends in a merge (check mark), a discard (gray cross, with its
reason) or an open head on NOW, continued as a dashed line into the gate of
its planned target. Glyphs for several lines or streams are shown in
\cref{fig:hardware-fork,fig:backport,fig:shared} and in \cref{app:card}.}
\label{fig:anatomy}
\end{figure}


\paragraph{Four kinds of branch.}
Every branch has exactly one type, shown by its lane, its color and the
letter in its start node: \textbf{I}~improvement (a change that makes the
product better for users), \textbf{L}~laboratory evolution (evaluation of
a new component, design or technology), \textbf{X}~issue (a defect) and
\textbf{T}~task (engineering or operational work that keeps the product
correct and operable). The split is deliberate: I and L move the product
forward and sit above the line; X and T keep it healthy and sit below.

\paragraph{The seven rules.}
HPB is defined by seven rules. They are short because each one closes a
specific hole I kept falling into. Their formal statements are in
\cref{app:formal}; here is what they say and why they exist.
\begin{description}[leftmargin=1.2em,style=nextline,itemsep=3pt,topsep=4pt]
  \item[\Cref{rule:origin}: single origin.]
    Every branch forks from exactly one released baseline, fixed when the
    work is accepted. For a defect, that is the version where it was seen.
    \emph{Why:} without it, ``relative to which version?'' has no answer
    and you cannot tell in which version a defect escaped.
  \item[\Cref{rule:closure}: explicit closure.]
    Every closed branch ends in exactly one way: merged into one baseline,
    or discarded with a written reason. Nothing is closed silently.
    \emph{Why:} ``closed'' with no version and no reason is the most common
    way traceability is lost. Six months later nobody can tell whether it
    was fixed, dropped or forgotten.
  \item[\Cref{rule:time}: time moves forward.]
    A branch merges only into a baseline that comes after its origin, on
    the same line or a descendant of it. \emph{Why:} a version cannot
    contain work that was requested after it was frozen, so change logs
    contain only what really went in.
  \item[\Cref{rule:gate}: gate integrity.]
    A planned version is released only when every branch committed to it
    has passed verification against a candidate of that version, and
    everything it depends on is released too. Work that is not ready is
    re-targeted or discarded, explicitly. \emph{Why:} it stops the most
    expensive habit in product work, shipping a version while committed
    work is quietly still pending.
  \item[\Cref{rule:lineage}: lineage and life cycle.]
    A hardware change that is not interchangeable in form, fit or function
    creates a new line with a permanent link to its parent; references are
    never edited in place, and a retired reference does not accumulate new
    work. \emph{Why:} units already in the field keep their history, and
    you always know what to maintain and for how long.
  \item[\Cref{rule:typed}: typed work.]
    Every branch has one type; sub-tasks and test runs are steps of a
    branch, not branches of their own. \emph{Why:} otherwise the same piece
    of work is counted twice, and the metrics lie.
  \item[\Cref{rule:discussion}: traceable discussion.]
    Every comment, decision or piece of evidence is attached to one
    element of the graph (a line, a baseline or a branch), and the record
    is append-only. \emph{Why:} decisions made in chat threads and meetings
    cannot be audited; attached to the branch they concern, they are found
    exactly where the diagram is read.
\end{description}

\paragraph{HPB 1.0: what is normative.}
This section, together with the quick-reference card in \cref{app:card},
is the normative definition of HPB~1.0: the idea, the four branch types,
the seven rules and the notation. Everything else in this paper explains,
illustrates or supports it, and the formal model of \cref{app:formal}
states the same rules precisely. The method will change only through
versioned releases: HPB~1.x for additions that keep every existing diagram
valid, a new major version for anything that would not. The text of the
method and its notation are offered under the Creative Commons Attribution
4.0 license (CC~BY~4.0): copy, translate and adapt them, with credit. The
\LaTeX{} package \texttt{hpb.sty} and the HPB Markup toolkit are released
under the MIT license.

\paragraph{What you get.}
With these rules in place, the questions from \cref{sec:introduction} are
answered by looking at the picture: where a defect was seen (its fork),
whether it is fixed and in which version (its terminal), what else is open
on the product (the branches crossing NOW), what a planned version is
waiting for (the branches entering its gate), and what moved to a new
hardware reference (the line fork). The rest of the paper makes each of
these elements precise (\cref{sec:model}), fixes how they are drawn
(\cref{sec:notation}), and shows them in the situations a product manager
meets every week (\cref{sec:practice}).
